\section{Hilbert--valued one--M\"obius compression}
\label{sec:tpc41-hilbert}

The four--factor expansion is exact but can obscure the operator that
must be estimated.  We now compress each fixed--alias column whose norm
square generates the corresponding fixed--alias row form: one target
M\"obius factor becomes the outer scalar, while all row information is
retained in a Hilbert coefficient.

\subsection{An exact vector identity}

Fix an alias \(k\).  Let
\begin{equation}
 a_k(r)=\log(r+h+kM)\Phi_0(r+h+kM),
 \label{eq:tpc41-hilbert-ak}
\end{equation}
interpreted as zero outside the positive terminal support.  Define a
vector \(\mathbf G_k^L(r)\in\ell^2\{(\gamma,j)\}\) by
\begin{equation}
 [\mathbf G_k^L(r)]_{\gamma,j}
 =\sum_{\alpha:m_\alpha j=r}
   c^L_{\alpha,\gamma,j}a_k(r).
 \label{eq:tpc41-hilbert-G}
\end{equation}
For fixed \((\gamma,j,r)\), source separation leaves at most one
\(\alpha\) in this sum.

For comparison with the inherited atomic diagonal, set
\begin{equation}
 \cD_{k,L}
 =\sum_{\gamma,j,\alpha}
 |c^L_{\alpha,\gamma,j}|^2
 |\mu(m_\alpha j+h+kM)|^2|a_k(m_\alpha j)|^2.
 \label{eq:tpc41-Hilbert-actual-diagonal}
\end{equation}

\begin{proposition}[Exact Hilbert compression]
\label{prop:tpc41-Hilbert-compression}
The fixed--alias terminal column satisfies
\begin{equation}
 \boxed{
 \mathscr T_k^L
 =-\sum_r\mu(r+h+kM)\mathbf G_k^L(r).}
 \label{eq:tpc41-Hilbert-compression}
\end{equation}
Moreover, its deterministic coefficient energy is
\begin{equation}
 \sum_r\|\mathbf G_k^L(r)\|_2^2
 =\sum_{\gamma,j,\alpha}
 |c^L_{\alpha,\gamma,j}|^2
 |a_k(m_\alpha j)|^2
 =:\widetilde{\cD}_{k,L}.
 \label{eq:tpc41-Hilbert-coefficient-energy}
\end{equation}
The actual atomic diagonal \(\cD_{k,L}\), which includes
\(|\mu(r+h+kM)|^2\), is at most \(\widetilde{\cD}_{k,L}\); both obey the
same deterministic endpoint upper bound.  The right column has the
identical formulation.
\end{proposition}

\begin{proof}
Insert \(r=m_\alpha j\) into
\eqref{eq:tpc41-interface-left-output} and group terms with the same
\(r\).  This gives \eqref{eq:tpc41-Hilbert-compression}.  When the norm in
\eqref{eq:tpc41-hilbert-G} is squared, different \(j\) are different
coordinates, and source separation prevents two rows from contributing to
the same fixed \((j,r)\).  Hence no cross term appears and
\eqref{eq:tpc41-Hilbert-coefficient-energy} follows.
\end{proof}

The equality--alias target can now be stated as a vector M\"obius
disjointness inequality.  Since
\(\widetilde{\cD}_{0,L}+\widetilde{\cD}_{0,R}\ll X^{o(1)}Q^2J\), a
sufficient estimate is
\begin{equation}
 \left\|\sum_r\mu(r+h)\mathbf G_0^L(r)\right\|_2^2
 +\left\|\sum_r\mu(r+h)\mathbf G_0^R(r)\right\|_2^2
 \ll X^{o(1)}K_B
 \sum_r\bigl(
 \|\mathbf G_0^L(r)\|_2^2+
 \|\mathbf G_0^R(r)\|_2^2\bigr).
 \label{eq:tpc41-Hilbert-equality-target}
\end{equation}
The coefficient--blind constant is of order \(Q\); the target constant is
only \(K_B\), so the required aggregate row gain is
\(Q/K_B=X^{o(1)}J^2\).

For the full alias trace, a cleaner but stronger sufficient theorem is
\begin{align}
 &\sum_{k=-L}^{L}
 \left\|\sum_r\mu(r+h+kM)\mathbf G_k^L(r)\right\|_2^2
 \notag\\
 &\quad+
 \sum_{k=-L}^{L}
 \left\|\sum_r\mu(r+h+kM)\mathbf G_k^R(r)\right\|_2^2
 \notag\\
 &\qquad\ll X^{o(1)}
 \sum_{k,r}\bigl(
 \|\mathbf G_k^L(r)\|_2^2+
 \|\mathbf G_k^R(r)\|_2^2\bigr).
 \label{eq:tpc41-Hilbert-trace-target}
\end{align}
This asks for constant--scale fixed--alias row decorrelation after the
full alias trace has deleted every cross--alias product.

\subsection{Why scalar compression is insufficient}

If the physical output coordinates are suppressed, one may formally
write a row sum as
\begin{equation}
 R_h
 =\sum_{\ell,d,j}
   \mu(d)(\log\ell)W(\ell,d,j)
   \mu(\ell dj+h)
 =\sum_nG(n)\mu(n+h),
 \label{eq:tpc41-scalar-compression}
\end{equation}
where
\begin{equation}
 G(n)=\sum_{\ell dj=n}\mu(d)(\log\ell)W(\ell,d,j).
 \label{eq:tpc41-scalar-G}
\end{equation}
This is an identity, not a cancellation estimate.  In particular:
\begin{enumerate}[label=(\roman*)]
 \item the divisor M\"obius factor is stored inside \(G\), not removed;
 \item the physical quantity is a sum of squared coordinate outputs,
 not the square of a scalar sum over coordinates; and
 \item dualizing the Hilbert norm introduces arbitrary \(\ell^2\) dual
 coefficients, so a scalar estimate that cancels between \((\gamma,j)\)
 coordinates does not imply \eqref{eq:tpc41-Hilbert-equality-target}.
\end{enumerate}

TPC--36 nevertheless supplies meaningful structure.  For each fixed
opposite row \(\gamma\), its projective reassembly decomposes the actual
coefficient into signed layers of total mass \(X^{o(1)}\), each of the
form
\begin{equation}
 b_\gamma(\ell)f_\gamma(d)\eta_\gamma(j/J),
 \label{eq:tpc41-projective-layer}
\end{equation}
with the \(\ell\)-factor independent of \((d,j)\).  Thus a uniform
Hilbert--valued M\"obius theorem for this restricted layer class could be
reassembled by Minkowski with only an \(X^{o(1)}\) loss.  The existing
projective decomposition does not itself prove that theorem: the common
opposite row, gcd mask, row separation, residue cell, and orbit weight
remain part of the vector coefficient.

\begin{assumption}[Restricted vector M\"obius input]
\label{ass:tpc41-restricted-vector-input}
For every projective layer generated by TPC--36, uniformly in its
opposite row and admissible alias, the corresponding vectors satisfy
either \eqref{eq:tpc41-Hilbert-equality-target} or the stronger trace
estimate \eqref{eq:tpc41-Hilbert-trace-target}.
\end{assumption}

\Cref{ass:tpc41-restricted-vector-input} is a precise continuation
interface, not a result of the present paper.  It is coefficient--specific
and therefore lies outside the universal positive phase methods ruled out
by TPC--40.
